\section{Boxes and their measures}
\subsection{Right-handed boxes}
\label{sec:boxes}

We now turn to the outline of the construction and its principal
ingredients.
Let $f\colon \mathbb{R}\Curves^+(\Sigma) \to \mathbb{R}$ be a
functional satisfying symmetry, additive union, smoothing, stability
and homogeneity.
Recall that since $f$ satisfies the symmetry condition,
$f([a])=f([a^{-1}])$ for any $a \in \pi_1(S,\ast)$.
Recall, also, that the fact that $f$ satisfies the smoothing condition
for unoriented curves means $f$ is non-increasing under smoothing
essential crossings for either of the two possible unoriented resolutions
(the one coherent with the chosen orientation, and the one which
isn't).

We wish to construct a current $\mu_f$ and show that
$f(C) = i(C, \mu_f)$. To construct~$\mu_f$, as mentioned above, we fix
(arbitrarily) a simple closed oriented curve~$[x]$ passing through the
basepoint. The endpoints of the lifts of $[x]$ are dense in
$\partial_{\infty} \Sigma$, by minimality of the action of $\pi_1(S)$ (Lemma~\ref{lem:minimal}).
% \DMGtodo[inline]{(Punctures) For surface with punctures, it is density in the limit set, not the whole circle.}

The main objects in our proof are certain boxes in the
space of geodesics to which we will assign measures. The prototype for
the kind of boxes we consider is shown in
Figure~\ref{fig:newmodel3}: we have four lines, all lifts of~$x$, with
two lifts $\ell_1,\ell_2$ on the left and two $r_1,r_2$ on the right,
and consider the window
$[\ell_2^-,\ell_1^-) \times [r_2^+,r_1^+) \subset
\partial_\infty \wt S \times \partial_\infty \wt S$ as a basic domain
for defining the measure. The lifts are defined by a 4-tuple of
elements $[p_1,q_1,p_2,q_2]$ of $\pi_1(S)$ by
$\ell_i = (p_i)^{-1} \cdot \wt x$ and $r_i = q_i \cdot \wt x$, with
conditions explained below.

\begin{figure}
\centering{
\resizebox{100mm}{!}{\svgfig{newmodel3}}
\caption{Schematically determining a box from group theory data.}
\label{fig:newmodel3}
}
\end{figure}

The measure that we ultimately assign to the associated window should
be invariant under
\begin{itemize}
\item changing a $p_i$ preserving $\ell_i$, by $p_i \mapsto xp_i$;
\item changing a $q_i$ preserving $r_i$, by $q_i \mapsto q_ix$; and
\item translating all the lifts by $g \in \pi_1(S)$:
  \[
    [p_1,q_1,p_2,q_2] \mapsto g \cdot [p_1,q_1,p_2,q_2] \coloneqq
      [p_1 g^{-1}, gq_1, p_2 g^{-1},  gq_2].
  \]
\end{itemize}
In fact, ultimately we will take a limit of applying the first
transformations: we look at the limit as $n\to \infty$ of $p_i \mapsto x^np_i$
and $q_i \mapsto q_i x^n$.

  If we are only interested in the relative position of these lines
  (for instance to construct $\pi_1(S)$-invariant functionals), we
  can parameterize them instead by group elements $(a,b,c,d)$ related
  to the $p_i$ and $q_i$ by
\begin{equation}
\begin{aligned}
  a &= p_1 q_1 &\qquad\qquad\qquad\qquad\qquad b &= p_1 q_2\\[2pt]
  c &= p_2 q_1 & d &= p_2 q_2.
\end{aligned}
  \label{eq:relationdatum}
\end{equation}
\DPTtodo[inline]{Necessarily have $aCdB = 1$}
Note that, by Equation~\eqref{eq:double_coset}, the relative position
of (say) $\ell_1$ and $r_1$ is given by the double coset
$\langle x \rangle a \langle x \rangle$.

The action of translating all lifts by $\pi_1(S)$ has been baked in to
this alternate parametrization,
but the measure corresponding to the associated window should be
invariant under the following four transformations:
\begin{itemize}
\item translating on $\ell_1$: $a \mapsto xa$ and $b \mapsto xb$;
\item translating on $\ell_2$: $c \mapsto xc$ and $d \mapsto xd$;
\item translating on $r_1$: $a \mapsto ax$ and $c \mapsto cx$; and
\item translating on $r_2$: $b \mapsto bx$ and $d \mapsto dx$.
\end{itemize}

Conversely, given elements $a,b,c,d \in \pi_1(S)$ with
$ac^{-1}db^{-1} = 1$, we can
construct four lifts of $x$ with the above relations, for instance
setting $p_1 = 1$, $q_1 = a$, $q_2 = b$, and $p_2 = ca^{-1}$.

As mentioned above, we need some conditions on these lifts; for
instance, we need the window
$[\ell_2^-,\ell_1^-) \times [r_2^+,r_1^+)$ to avoid the
diagonal. For the proofs to go through easily, we will impose more
conditions on the configurations of these lifts.

\begin{definition}[RHS lift system]
\label{def:RH-lift-sys}
Let $a,b,c,d \in \pi_1(S,\ast)$ so that:
\begin{enumerate}
\item\label{item:box-close} $a c^{-1} d b^{-1} = 1$;
\item\label{item:box-edges-cross} $a,b,c,d$ cross $x$ to the right; and
\item\label{item:box-diags-cross} for sufficiently large~$n$,
  $dx^na c^{-1}
  = cAbx^na c^{-1}$ crosses
  $cx^n$ to the right.
\end{enumerate}
Then we call the datum $(a,b,c,d)$ a \emph{right-handed lift system}.
As explained above, this determines a $\pi_1(S)$-orbit of windows. In
terms of the alternative data $[p_1,q_1,p_2,q_2]$ for a window (rather
than for an orbit), condition~(1) is automatic and the others
become
\begin{enumerate}[label=(\arabic*$^\prime$),start=2]
\item for $i,j \in \{1,2\}$, $p_iq_j$ crosses $x$ to the right; and
\item for sufficiently large~$n$, $q_2x^np_1$ crosses $q_1x^np_2$ to
  the right.
\end{enumerate}
\end{definition}
See Figure~\ref{fig:box4} for examples of
RH lift systems, and Figure~\ref{fig:box4} for non-examples.
 We
distinguish three types of RH lift systems: type 1 has
no nested lifts; type 2 has one parallel pair of nested lifts; and
type 3 has two parallel pairs of nested lifts of $x$.
Note that since $x$ is simple, lifts
cannot intersect each other, and nested lifts cannot be anti-parallel
by the right-handed condition.


\begin{figure}[h!]
     \centering
     \begin{subfigure}[b]{0.45\textwidth}
\centering{
\resizebox{30mm}{!}{\svgfig{box1}}}
     \end{subfigure}
     \hfill
     \begin{subfigure}[b]{0.45\textwidth}
\centering{
\resizebox{30mm}{!}{\svgfig{box}}}
     \end{subfigure}
          \begin{subfigure}[b]{0.45\textwidth}
\centering{
\resizebox{30mm}{!}{\svgfig{box4}}}
     \end{subfigure}
     \caption{Examples of right-handed boxes}
     \label{fig:rhbox}
\end{figure}

\begin{remark}
%\begin{enumerate} 
%\item
The reader might wonder why in condition~\ref{item:box-diags-cross} we require the condition
  ``for $n$ large enough''. It turns out that requiring the condition
  for all $n$ is too restrictive, and while it does happen for some
  configurations of lifts, such as those of Type 1,
  in general it doesn't hold for the types of configurations of lifts
  that we are interested in. 
 Consider Figure~\ref{fig:nestedpositivity}: item a) shows, $[bx^n]$ and $[cx^n]$ might not cross. However, for sufficiently large $n$, the $[b x^n]$ and $[c x^n]$ will cross, as item b) suggests. To prove it, consider the bands model centered at the outer most nested lift of $x$, as in item c). For sufficiently large $n$, it's clear that $[b x^n]$ and $[c x^n]$ cross each other.
 There are some configurations of boxes, such as Type 1, Figure~\ref{fig:rhbox}, where the crossing is essential without passing to large $n$.
%\item Condition (3) in the definition of RH lift is equivalent to the following statement: there exists $n_0$ so that for all $n \geq n_0$, the lift of the path $ca^{-1}ax^n a c^{-1}$ and the path $dx^n$, both based at $*$, are not linked.
%$$  \DPTtodo[inline]{This is nice. Probably worth spelling out some of the other
%    implications about the order in which things appear on the
%    boundary. But is the other direction correct? Could we end up
%    with a left-handed crossing?}
 % To see this, note that, for any $n$, the startpoint $(dx^n)^-$ appears before $\lim_n (dx^n)^-$ in counter-clockwise order. By item (3), $\lim_n (ca^{-1}ax^n a c^{-1})^-$ appears before $\lim_n (dx^n)^-$, and so for large enough $n$, $(dx^n)^-$ appears before $(ca^{-1}ax^n a c^{-1})^-$. The other implication is similar.
%     \end{enumerate}
\label{rmk:largeenough}
\end{remark}

\begin{definition}\label{def:RH-box}
Fix once and for all a lift $\tilde{x}$ of $[x]$. Every right handed lift system given by the datum $[p_1,p_2,q_1,q_2]$  determines uniquely a \emph{right-handed box}, a half-closed half-open box in $\partial_{\infty} \Sigma \times \partial_{\infty} \Sigma - \Delta$ given by, following the notation above,
\[
B_{p_1,p_2,q_1,q_2} \coloneqq [\ell_2^-,\ell_1^-) \times [r_2^+,r_1^+)
\]

By construction, a different choice of lift $\tilde{x}$ determines a different box $B_{p_1,p_2,q_1,q_2}'$, whose endpoints differ from the endpoints of $B_{p_1,p_2,q_1,q_2}$ by translation by some element in $\pi_1(S)$.
On the other hand, the RH lift system $(a,b,c,d)$
determines a $\pi_1(S)$-orbit of boxes that we will denote
$[B_{p_1,p_2,q_1,q_2}]$ or $[B]_{a,b,c,d}$ depending on if we want to
draw attention to a concrete representative of the orbit or to the
orbit itself. 

We will denote by $\mathcal{B}$ the set of right-handed boxes.
\end{definition}


Given a right-handed lift system $(a,b,c,d)$ and a choice of lift $\tilde{x}$,  this yields configurations of lifts of Type 1, Type 2 and Type 3, where the lifts are placed as shown in Figure~\ref{fig:rhbox}. respectively.
Note how the lifts are non-intersecting, since $[x]$ is simple, and these orientation of the lifts are always right-handed.
\begin{example}
The box in Figure~\ref{fig:box4} is \emph{not} a right-handed box, since $c$ crosses a lift of $x$ left-handedly, and $d$ crosses a lift of $x$ left-handedly.
\begin{figure}
\centering{
\resizebox{60mm}{!}{\svgfig{nonbox}}
\resizebox{60mm}{!}{\svgfig{box3}}
\caption{Two non-examples of right-handed boxes. The left is not
  right-handed because $b$ and $c$ cross in the wrong way, while the
  right is a non-example because $c$ does not cross $x$ essentially(to
  the right
  or to the left)}
\label{fig:box4}
}
\end{figure}
\end{example}

\begin{remark}
  The reader may wonder why we choose to make the boxes half-open in
  the specified way. In fact it doesn't matter, as by
  Proposition~\ref{prop:spikes}
  the weight of the
  boundary of these boxes under any geodesic current is necessarily
  zero: the only closed geodesic with an endpoint starting at the base
  of any of the involved spikes is a lift of~$x$, and by definition of
  right-handed box, it is not contained in the spikes. This choice is
  more convenient for the proofs.
\end{remark}

In our construction, we want to use the following fact.

\begin{lemma}
The action of $\pi_1(S)$ on $\partial_{\infty} \tilde{S}$ is minimal, i.e., for any $p \in \partial_{\infty} \tilde{S}$, the orbit $\pi_1(S)p$ is dense.
\label{lem:minimal}
\end{lemma}
% \DMGtodo[inline]{(Punctures) Again, this uses surface is closed.}
However, because of our interest in right-handed boxes, we are only considering a subset of the lifts of $[x]$ to construct our boxes.
These are still enough to approximate any box.
% \DMGtodo[inline]{(Punctures) This uses surface is closed (commented before, too)}



\begin{lemma}
Let $x \in \pi_1(\Sigma)$.
Given any $z \in \partial \Sigma$,
% \DMGtodo[inline]{(Punctures) Uses
%   closed surface, otherwise replace by limit set.}
and any $\epsilon>0$, there exists a lift $g \cdot \wt{x}$ whose endpoints $g\wt{x}_+$ and $g\wt{x}_-$ are in $(z-\epsilon,z+\epsilon)$, and they are ordered counter-clockwise (or clockwise)
with respect to the $ \partial \Sigma$ order.
\label{lem:diagonaldense}
\end{lemma}
\begin{proof}
By Lemma~\ref{lem:minimal}, endpoints of axes of $\pi_1(\Sigma)$ elements are dense in $\partial \Sigma$. Consider an element $g$ so that $\wt{x}$ is parallel (resp. anti-parallel) to $\ora{g}$ and $g^+ \in (z-\epsilon,z+\epsilon)$ (see Figure~\ref{fig:approx2}). Then, for large enough $n$, by north-south dynamics of $g$, we have that $(g^n x g^{-n})^{\pm} \in (z-\epsilon,z+\epsilon)$ in clockwise (resp. counter-clockwise) order.
\end{proof}

\begin{proposition}
The family of right handed boxes $\mathcal{B}$ is dense in the space
of rectangles in $G^+(\Sigma)$, in the following sense.
For any rectangle $R = [p,q) \times [r,s) \subset G^+(\Sigma)$,
and any $\epsilon>0$, there exist $a,b,c,d \in \pi_1(S)$ forming a
right-handed box so that for some $B \in [B]_{a,b,c,d}$, with
$B=[p',q') \times [r',s')$, we have $p' \in B_{\epsilon}(p)$, $q' \in
B_{\epsilon}(q)$, $r' \in B_{\epsilon}(r)$ and $s' \in
B_{\epsilon}(s)$. 
\label{prop:density}
\end{proposition}
\begin{proof}\DMGtodo[inline]{do we need to explain why all the conditions of right handed box are met? }
   Consider a lift $\tilde{x}$ of $[x]$ with startpoint $x^-$ and endpoint $x^+$.
   Given a points $p,q,r,s \in \partial_{\infty} \tilde{\Sigma}$ in this clockwise order, and any $\epsilon>0$, we will produce lifts $\tilde{x}_p,\tilde{x}_q,\tilde{x}_r,\tilde{x}_s$ of $[x]$ with endpoints within $\epsilon$ of $p,q,r,s$, respectively, and with the right orientations.
   For $r$ (resp. $s$), by Lemma~\ref{lem:diagonaldense}, we find a translate $g_r$ of $\wt{x}$ with endpoints in $B_{\epsilon}(r)$ (resp. a translate $g_s$ with endpoints in $B_{\epsilon}(r)$), with clockwise orientations.


\begin{figure}
\centering{
\resizebox{60mm}{!}{\svgfig{approx2}}
\resizebox{60mm}{!}{\svgfig{approxfin}}
\caption{Approximating a boundary point of the interval of endpoints of a right-handed box}
\label{fig:approx2}
}
\end{figure}

   For $p$ (resp. $q$), by Lemma~\ref{lem:diagonaldense}, we find a translate $g_p$ of $\wt{x}$ with endpoints in $B_{\epsilon}(p)$ (resp. a translate $g_q$ with endpoints in $B_{\epsilon}(q)$), with counter-clockwise orientations. See right Figure~\ref{fig:approx2}.
\end{proof}

\subsection{Measures of boxes}

Fix $x$ and a lift $\tilde{x}$ once and for all.

\begin{proposition}
\label{prop:limitexists}
Let $B_{p_1,p_2,q_1,q_2}$ be a right-handed box so that $[B_{p_1,p_2,q_1,q_2}]=[B]_{a,b,c,d}$.
The following limit exists:
\[
\mu_{f}(B_{p_1,p_2,q_1,q_2})\coloneqq\frac{1}{2}\lim_{n \to \infty}f([bx^n]) + f([cx^n])-f([ax^n]) - f([dx^n]).
\]
\end{proposition}
Note this function on boxes (which will eventually become a pre-measure and
then a measure) is automatically $\pi_1(S)$-invariant.

We prove Proposition~\ref{prop:limitexists} by breaking it into several lemmas.





%%% Local Variables:
%%% mode: latex
%%% TeX-master: "Intersections"
%%% End:
